% from ega1-10-fr.tex
\oldpage[I]{191}

soit de type fini sur
$\mathscr{O}_0=\mathscr{O}_1/\mathscr{J}_1$. Alors $\mathfrak{X}$ est
un préschéma formel adique, et si $\mathscr{J}^{(n)}$ est le noyau de
$\mathscr{O}_{\mathfrak{X}}\to\mathscr{O}_n$, on a
$\mathscr{J}^{(n)}=\mathscr{J}^{n+1}$ et
$\mathscr{J}/\mathscr{J}^2$ est isomorphe à $\mathscr{J}_1$. Si en
outre $X_0$ est localement noethérien (resp. noethérien),
$\mathfrak{X}$ est localement noethérien (resp. noethérien).
\end{corollary}

\begin{proof}
Comme les espaces sous-jacents à $\mathfrak{X}$ et $X_0$ sont les mêmes,
la question est locale et l'on peut supposer tous les $X_i$ affines~;
compte tenu des relations
$\mathscr{J}_{ij}=\widetilde{\mathfrak{J}_{ji}}$ (avec les notations
de (\hyperref[I.10.6.3-fr]{10.6.3})), on est aussitôt ramené aux
assertions correspondantes de
(\hyperref[0.7.2.7-fr]{0, 7.2.7} et
\hyperref[0.7.2.8-fr]{7.2.8}), en notant que
$\mathfrak{J}_1/\mathfrak{J}_1^2$ est alors un $A_0$-module de type
fini (\hyperref[I.1.3.9-fr]{1.3.9}).
\end{proof}

En particulier, \emph{tout préschéma formel localement noethérien
$\mathfrak{X}$} est limite inductive d'une suite $(X_n)$ de préschémas
(usuels) localement noethériens vérifiant les conditions de
(\hyperref[I.10.6.3-fr]{10.6.3}) et
(\hyperref[I.10.6.4-fr]{10.6.4})~: il suffit de considérer un Idéal de
définition $\mathscr{J}$ de $\mathfrak{X}$
(\hyperref[I.10.5.4-fr]{10.5.4}) et de prendre
$X_n=(\mathfrak{X},
\mathscr{O}_{\mathfrak{X}}/\mathscr{J}^{n+1})$
((\hyperref[I.10.5.1-fr]{10.5.1}) et
(\hyperref[I.10.6.2-fr]{10.6.2})).

\begin{corollary}[10.6.5]
\label{I.10.6.5-fr}
Soit $A$ un anneau admissible. Pour que le schéma formel affine
$\mathfrak{X}=\operatorname{Spf}(A)$ soit noethérien, il faut et il
suffit que $A$ soit adique et noethérien.
\end{corollary}

\begin{proof}
La condition est évidemment suffisante. Inversement, supposons que
$\mathfrak{X}$ soit noethérien, et soient $\mathfrak{J}$ un idéal de
définition de $A$, $\mathscr{J}=\mathfrak{J}^\Delta$ l'Idéal de
définition correspondant de $\mathfrak{X}$. Les préschémas (usuels)
$X_n=(\mathfrak{X},
\mathscr{O}_{\mathfrak{X}}/\mathscr{J}^{n+1})$ sont alors affines et
noethériens, donc les anneaux
$A_n=A/\mathfrak{J}^{n+1}$ sont noethériens
(\hyperref[I.6.1.3-fr]{6.1.3}), d'où on conclut que
$\mathfrak{J}/\mathfrak{J}^2$ est un $A/\mathfrak{J}$-module de type
fini. Comme les $\mathscr{J}^n$ forment un système fondamental d'Idéaux
de définition de $\mathfrak{X}$
(\hyperref[I.10.5.1-fr]{10.5.1}), on a
$\mathscr{O}_{\mathfrak{X}}
=\varprojlim(\mathscr{O}_{\mathfrak{X}}/\mathscr{J}^n)$
(\hyperref[I.10.5.3-fr]{10.5.3})~; on en conclut
(\hyperref[I.10.1.3-fr]{10.1.3}) que $A$ est topologiquement isomorphe à
$\varprojlim A/\mathfrak{J}^n$, donc est adique et noethérien
(\hyperref[0.7.2.8-fr]{0, 7.2.8}).
\end{proof}

\begin{remark}[10.6.6]
\label{I.10.6.6-fr}
Avec les notations de (\hyperref[I.10.6.3-fr]{10.6.3}), soit
$\mathscr{F}_i$ un $\mathscr{O}_i$-Module, et supposons donné, pour
$i\geq j$, un $v_{ij}$-morphisme
$\theta_{ji}:\mathscr{F}_i\to\mathscr{F}_j$, de sorte que
$\theta_{kj}\circ\theta_{ji}=\theta_{ki}$ pour $k\leq j\leq i$.
Comme l'application continue sous-jacente à $v_{ij}$ est l'identité,
$\theta_{ji}$ est un homomorphisme de faisceaux de groupes abéliens sur
l'espace $\mathfrak{X}$~; en outre, si $\mathscr{F}$ est la limite
projective du système projectif $(\mathscr{F}_i)$ de faisceaux de groupes
abéliens, le fait que les $\theta_{ji}$ sont des $v_{ij}$-morphismes
permet de définir sur $\mathscr{F}$ une structure de
$\mathscr{O}_{\mathfrak{X}}$-Module par passage à la limite projective~;
muni de cette structure, nous dirons que $\mathscr{F}$ est la
\emph{limite projective} (pour les $\theta_{ji}$) du système de
$\mathscr{O}_i$-Modules $(\mathscr{F}_i)$. Dans le cas particulier où
$v_{ij}^*(\mathscr{F}_i)=\mathscr{F}_j$ et où $\theta_{ji}$ est
l'\emph{identité}, nous dirons pour abréger, que $\mathscr{F}$ est la
limite projective d'un système $(\mathscr{F}_i)$ tel que
$v_{ij}^*(\mathscr{F}_i)=\mathscr{F}_j$ pour $j\leq i$ (sans
mentionner les $\theta_{ji}$).
\end{remark}

\begin{env}[10.6.7]
\label{I.10.6.7-fr}
Soient $\mathfrak{X}$, $\mathfrak{Y}$ deux préschémas formels,
$\mathscr{J}$ (resp. $\mathscr{K}$) un Idéal de définition de
$\mathfrak{X}$ (resp. $\mathfrak{Y}$),
$f:\mathfrak{X}\to\mathfrak{Y}$ un morphisme tel que
$f^*(\mathscr{K})\mathscr{O}_{\mathfrak{X}}\subset\mathscr{J}$. On a
alors pour tout entier $n>0$,
$f^*(\mathscr{K}^n)\mathscr{O}_{\mathfrak{X}}
=(f^*(\mathscr{K})\mathscr{O}_{\mathfrak{X}})^n
\subset\mathscr{J}^n$~; on peut donc
(\hyperref[I.10.5.6-fr]{10.5.6}) déduire de $f$ un morphisme de
préschémas (usuels) $f_n:X_n\to Y_n$, en posant
$X_n=(\mathfrak{X},
\mathscr{O}_{\mathfrak{X}}/\mathscr{J}^{n+1})$,
$Y_n=(\mathfrak{Y},
\mathscr{O}_{\mathfrak{Y}}/\mathscr{K}^{n+1})$, et il résulte
aussitôt des définitions que les diagrammes
\[
  \label{I.10.6.7.1-fr}
  \xymatrix{
    X_m\ar[r]^{f_m}\ar[d] &
    Y_m\ar[d]\\
    X_n\ar[r]_{f_n} &
    Y_n
  }
  \tag{10.6.7.1}
\]

